\documentclass[11pt,a4paper]{article}
\usepackage[utf8]{inputenc}
\usepackage[russian]{babel}
\usepackage{amsmath, amssymb, amsfonts}
\usepackage{booktabs}
\usepackage{geometry}
\usepackage{array}
\usepackage{microtype}

\geometry{top=2cm, bottom=2cm, left=2cm, right=2cm}

\begin{document}

\section{Полное описание используемого метода}

Предложенный метод представляет собой легковесную и интерпретируемую систему автоматической классификации дефектов на картах кремниевых пластин (wafer maps) по 9 каноническим классам (\textit{Center, Donut, Edge-Loc, Edge-Ring, Loc, Random, Scratch, Near-Full, None}). Метод решает три основные проблемы: сильный дисбаланс классов, жесткие ограничения на вычислительные ресурсы (параметры/память) и необходимость интерпретируемости решений.

Метод состоит из трех основных этапов:
\begin{enumerate}
    \item \textbf{Препроцессинг и балансировка данных (на этапе обучения):}
    \begin{itemize}
        \item Выбираются карты пластин строго в исходном разрешении $26 \times 26$ без изменения размера, чтобы избежать артефактов интерполяции.
        \item Каждая карта кодируется в трехканальный one-hot тензор размерности $26 \times 26 \times 3$ (каналы: фон/отсутствие чипа, годный чип, дефектный чип).
        \item Борьба с дисбалансом классов ведется за счет комбинации уменьшения масштаба (\textit{downsampling}) мажоритарного класса \textit{None} и генеративного расширения (\textit{data augmentation}) миноритарных классов с использованием сверточного автокодировщика (CAE).
    \end{itemize}

    \item \textbf{Легковесная архитектура сверточной нейросети с модифицированным CBAM (Proposed CBAM-CNN):}
    \begin{itemize}
        \item Состоит из двух блоков «Свертка — ReLU — MaxPool», после каждого из которых интегрирован модифицированный модуль внимания CBAM (\textit{Convolutional Block Attention Module}). Модификации направлены на адаптацию к низкому разрешению $26 \times 26$ и стабилизацию обучения при сильном дисбалансе классов.
        \item \textbf{Канальное внимание (Channel Attention):} Использует общий двухслойный MLP, но с добавлением слоев пакетной нормализации (\textit{Batch Normalization}) после каждого аффинного преобразования для стабилизации активаций на малых пространственных сетках.
        \item \textbf{Пространственное внимание (Spatial Attention):} Заменяет стандартную одиночную свертку на параллельный мультимасштабный блок с ядрами свертки $3 \times 3$, $5 \times 5$ и $7 \times 7$, выходы которых объединяются операцией усреднения. Это позволяет одновременно улавливать как мелкие царапины, так и крупные кольцевые/периферийные дефекты.
        \item \textbf{Резидуальные связи (Residual Skips):} Внутри каждого CBAM-блока добавлен сквозной проброс градиента (\textit{skip connection}) для предотвращения затухания градиентов и переподавления признаков.
    \end{itemize}

    \item \textbf{Классификационная голова:}
    \begin{itemize}
        \item Тензор признаков после второго CBAM сжимается (\textit{Flatten}) и пропускается через скрытый полносвязный слой (128 нейронов, ReLU) на выходной слой с 9 нейронами и функцией активации Softmax. Общий размер модели составляет всего около 153 тыс. параметров ($\sim$0.6 МБ в формате FP32).
    \end{itemize}
\end{enumerate}

\section{Математические формулы метода}

\subsection{Препроцессинг и кодирование}
Исходный labeled корпус дефектов $D_0$ фильтруется до подмножества $D_{26}$ карт с разрешением $26 \times 26$:
\begin{equation}
D_{26} = \left\{ X \in D_0 \;\Big\vert\; X \in \{0, 1, 2\}^{26 \times 26} \right\}
\end{equation}
где значения $0, 1, 2$ соответствуют состояниям: фон, годен, брак.

Каждая карта $X$ преобразуется в тензор $T \in \{0,1\}^{26 \times 26 \times 3}$:
\begin{equation}
T(i, j, c) = \mathbf{1}[X(i, j) = c], \quad c \in \{0, 1, 2\}
\end{equation}

\subsection{Математика сверточного бэкбона}
Дискретная 2D-свертка на слое $l$ выражается формулой:
\begin{equation}
\left(X^{(l-1)} * W^{(l)}\right)_{(i,j,c)} = \sum_{u=0}^{k-1} \sum_{v=0}^{k-1} \sum_{\gamma=1}^{C^{(l-1)}} W^{(l)}_{(u,v,\gamma,c)} X^{(l-1)}_{(i+u-p_t, \; j+v-p_l, \; \gamma)}
\end{equation}

Активация признаков с использованием функции ReLU:
\begin{equation}
F^l = \phi\left(X^{(l-1)} * W^{(l)} + b^{(l)}\right), \quad \phi(x) = \max(0, x)
\end{equation}

Размерность карты признаков после свертки со страйдом $s$ и паддингом $p_t, p_b, p_l, p_r$:
\begin{equation}
H_l = \left\lfloor \frac{H_{(l-1)} + p_t + p_b - k}{s} \right\rfloor + 1
\end{equation}
\begin{equation}
W_l = \left\lfloor \frac{W_{(l-1)} + p_l + p_r - k}{s} \right\rfloor + 1
\end{equation}

Операция MaxPooling $2 \times 2$ со страйдом 2:
\begin{equation}
P^l(i, j, c) = \max_{0 \le u, v < 2} F^l(2i+u, \; 2j+v, \; c), \quad H_{pool}^l = \lfloor H_l / 2 \rfloor, \quad W_{pool}^l = \lfloor W_l / 2 \rfloor
\end{equation}

Рекурсия эффективного рецептивного поля слоя $l$ ($r_l$) со страйдом $j_l$:
\begin{equation}
r_l = r_{(l-1)} + (k_l - 1)j_{(l-1)}, \quad j_l = j_{(l-1)}s_l \quad (\text{при } r_0 = j_0 = 1)
\end{equation}

Для архитектуры CBAM-CNN прогрессия имеет вид:
\begin{itemize}
    \item $\text{conv1 } (k=4, s=1) \Rightarrow r_1=4, j_1=1$
    \item $\text{pool1 } (k=2, s=2) \Rightarrow r_2=5, j_2=2$
    \item $\text{conv2 } (k=3, s=1) \Rightarrow r_3=9, j_3=2$
    \item $\text{pool2 } (k=2, s=2) \Rightarrow r_4=11, j_4=4$
\end{itemize}
Каждый юнит итоговой карты признаков $4 \times 4$ покрывает рецептивную область размером $22 \times 22$ от исходного изображения $26 \times 26$.

\subsection{Модифицированный блок CBAM}
Общая последовательность применения масок внимания:
\begin{equation}
X' = M_c(P_l) \odot P_l, \quad X_l = M_s(X') \odot X'
\end{equation}

\subsubsection*{Канальное внимание (Channel Attention)}
Векторы глобального усреднения (GAP) и максимизации (GMP):
\begin{equation}
s_{avg} = \text{GAP}(F) \in \mathbb{R}^C, \quad s_{max} = \text{GMP}(F) \in \mathbb{R}^C
\end{equation}

Функция обработки векторов в общем MLP со слоями пакетной нормализации ($\text{BN}_1, \text{BN}_2$) и коэффициентом сжатия $r=8$:
\begin{equation}
g(s) = \text{BN}_2\left(W_2 \delta\left(\text{BN}_1\left(W_1 s\right)\right)\right), \quad W_1 \in \mathbb{R}^{\frac{C}{r} \times C}, \; W_2 \in \mathbb{R}^{C \times \frac{C}{r}}
\end{equation}
\begin{equation}
M_c(F) = \sigma\left(g(s_{avg}) + g(s_{max})\right) \in [0, 1]^C
\end{equation}
где $\delta$ — активация ReLU, а $\sigma$ — сигмоида. Выходной канальный тензор:
\begin{equation}
F' = M_c(F) \odot F
\end{equation}

\subsubsection*{Пространственное внимание (Spatial Attention)}
Входной тензор $F'$ сжимается по каналам:
\begin{equation}
Q = \left[\text{AvgPool}_c(F'), \; \text{MaxPool}_c(F')\right] \in \mathbb{R}^{H \times W \times 2}
\end{equation}

Применение мультимасштабных параллельных сверток $k \times k$ ($k \in \{3, 5, 7\}$) и их усреднение:
\begin{equation}
A_k = \text{Conv}_{k \times k}(Q) \in \mathbb{R}^{H \times W \times 1}, \quad k \in \{3, 5, 7\}
\end{equation}
\begin{equation}
M_s(F') = \sigma\left(\text{Fuse}(A_3, A_5, A_7)\right), \quad \text{Fuse} = \frac{A_3 + A_5 + A_7}{3}
\end{equation}

Выход пространственного внимания:
\begin{equation}
\tilde{F} = M_s(F') \odot F'
\end{equation}

\subsubsection*{Резидуальное сложение}
\begin{equation}
Y = \tilde{F} + F
\end{equation}

\subsection{Полносвязный классификатор и функция потерь}
Преобразование в вектор $x \in \mathbb{R}^{1152}$ и прохождение через скрытый и выходной слои:
\begin{equation}
h = \text{ReLU}(W_1 x + b_1) \in \mathbb{R}^{128}, \quad z = W_2 h + b_2 \in \mathbb{R}^9
\end{equation}
\begin{equation}
p_k = \frac{e^{z_k}}{\sum_{j=1}^9 e^{z_j}}
\end{equation}

Обучение оптимизирует кросс-энтропию:
\begin{equation}
\mathcal{L}_{\text{CCE}} = - \frac{1}{N} \sum_{i=1}^N \sum_{j=1}^9 y_{ij} \log \hat{y}_{ij}
\end{equation}

\subsection{Математика автокодировщика (CAE)}
Кодирование:
\begin{equation}
Z = f_{\text{encoder}}(X) = \sigma(W_e * X + b_e), \quad Z \in \mathbb{R}^{13 \times 13 \times 64}
\end{equation}

Декодирование:
\begin{equation}
\hat{X} = f_{\text{decoder}}(Z) = \sigma(W_d * Z + b_d)
\end{equation}

Минимизация среднеквадратичной ошибки (MSE) при обучении:
\begin{equation}
\mathcal{L}_{\text{MSE}} = \frac{1}{N} \sum_{i=1}^N \Vert X_i - \hat{X}_i \Vert_2^2
\end{equation}

Добавление шума в латентное пространство:
\begin{equation}
Z_{\text{noisy}} = Z + \epsilon, \quad \epsilon \sim \mathcal{N}(0, \sigma^2)
\end{equation}
\begin{equation}
\hat{X}_{\text{noisy}} = f_{\text{decoder}}(Z_{\text{noisy}})
\end{equation}

\section{Все результаты экспериментов}

\subsection{Распределение классов до и после балансировки (в обучающей выборке)}
Балансировка применялась только к тренировочной выборке для предотвращения утечки данных (валидационный и тестовый наборы оставались строго реальными):

\begin{table}[h!]
\centering
\caption{Распределение классов в обучающей выборке}
\label{tab:class_distribution}
\begin{tabular}{lcc}
\toprule
\textbf{Класс дефекта} & \textbf{Кол-во до балансировки (\%)} & \textbf{Кол-во после CAE и даунсэмплинга (\%)} \\
\midrule
Center & 90 (0.63\%) & 2160 (10.96\%) \\
Donut & 1 (0.01\%) & 2002 (10.16\%) \\
Edge-loc & 296 (2.06\%) & 2368 (12.02\%) \\
Edge-ring & 31 (0.22\%) & 2046 (10.38\%) \\
Loc & 297 (2.07\%) & 2376 (12.06\%) \\
Near-full & 16 (0.11\%) & 2032 (10.31\%) \\
Random & 74 (0.52\%) & 2146 (10.89\%) \\
Scratch & 72 (0.50\%) & 2088 (10.60\%) \\
None & 13,489 (93.90\%) & 2489 (12.63\%) \\
\midrule
\textbf{Всего} & \textbf{14,366 (100.00\%)} & \textbf{19,707 (100.00\%)} \\
\bottomrule
\end{tabular}
\end{table}

\noindent (Класс \textit{None} был сокращен на 81.56\% — с 13 489 до 2 489 образцов).

\subsection{Сравнение с базовыми моделями по классам}
В таблице \ref{tab:comparison_baseline} приведено сравнение поклассовой точности предложенного метода с популярными глубокими архитектурами на тестовом наборе $26 \times 26 \times 3$ (метрики: Accuracy / Precision / Recall):

\begin{table}[h!]
\centering
\caption{Сравнение поклассовых метрик (Acc / Prec / Rec)}
\label{tab:comparison_baseline}
\small
\begin{tabular}{lccccccccc}
\toprule
\textbf{Модель} & \textbf{Center} & \textbf{Donut} & \textbf{Edge-loc} & \textbf{Edge-ring} & \textbf{Loc} & \textbf{Scratch} & \textbf{Near-full} & \textbf{Random} & \textbf{None} \\
\midrule
\textbf{ShuffleNet-v2} & 0.952 & 0.965 & 0.976 & 0.958 & 0.965 & 0.984 & 0.985 & 0.969 & 0.968 \\
& 0.938 & 0.951 & 0.959 & 0.945 & 0.938 & 0.972 & 0.961 & 0.968 & 0.954 \\
& 0.949 & 0.962 & 0.966 & 0.951 & 0.958 & 0.977 & 0.963 & 0.975 & 0.962 \\
\midrule
\textbf{ResNet50} & 0.770 & 0.900 & 0.880 & 0.950 & 0.590 & 0.700 & 0.950 & 0.870 & 0.890 \\
& 0.790 & 0.860 & 0.860 & 1.000 & 0.650 & 0.700 & 0.900 & 0.890 & 0.910 \\
& 0.750 & 0.950 & 0.900 & 0.900 & 0.550 & 0.700 & 1.000 & 0.850 & 0.880 \\
\midrule
\textbf{EfficientNetB0} & 0.900 & 0.900 & 0.870 & 1.000 & 0.820 & 0.680 & 0.920 & 0.850 & 0.920 \\
& 0.900 & 0.900 & 0.900 & 0.940 & 1.000 & 0.620 & 0.900 & 0.830 & 0.940 \\
& 0.900 & 0.900 & 0.900 & 1.000 & 0.700 & 0.750 & 0.950 & 1.000 & 0.900 \\
\midrule
\textbf{VGG16} & 0.890 & 0.900 & 0.770 & 0.950 & 0.630 & 0.560 & 0.950 & 0.830 & 0.880 \\
& 0.940 & 0.860 & 0.740 & 1.000 & 0.620 & 0.580 & 0.900 & 0.940 & 0.920 \\
& 0.850 & 0.950 & 0.810 & 0.900 & 0.650 & 0.550 & 1.000 & 0.750 & 0.830 \\
\midrule
\textbf{DC-net} & 0.990 & 0.978 & 0.965 & 0.944 & 0.998 & 0.980 & 0.958 & 0.934 & 0.991 \\
& 0.994 & 0.978 & 0.972 & 0.939 & 0.988 & 0.934 & 0.965 & 0.957 & 0.997 \\
& 0.983 & 0.982 & 0.967 & 0.942 & 0.967 & 0.983 & 0.952 & 0.929 & 0.995 \\
\midrule
\textbf{Proposed} & \textbf{0.990} & \textbf{1.000} & \textbf{0.990} & \textbf{1.000} & \textbf{0.980} & \textbf{0.990} & \textbf{1.000} & \textbf{1.000} & \textbf{0.990} \\
\textbf{CBAM-CNN} & \textbf{0.990} & \textbf{1.000} & \textbf{1.000} & \textbf{0.990} & \textbf{0.990} & \textbf{1.000} & \textbf{1.000} & \textbf{1.000} & \textbf{0.990} \\
& \textbf{1.000} & \textbf{1.000} & \textbf{0.980} & \textbf{1.000} & \textbf{0.960} & \textbf{0.990} & \textbf{1.000} & \textbf{1.000} & \textbf{0.980} \\
\bottomrule
\end{tabular}
\end{table}

\noindent Предложенная модель показывает превосходные результаты в сложных классах \textit{Loc} и \textit{Scratch} (полнота 0.96 и 0.99 соответственно) по сравнению с тяжелыми базовыми сетями.

\subsection{Сводное сравнение эффективности и сложности моделей}
В таблице \ref{tab:complexity_comparison} приведено сопоставление точности классификации, числа параметров и вычислительной сложности в GFLOPs:

\begin{table}[h!]
\centering
\caption{Сводные показатели точности и сложности}
\label{tab:complexity_comparison}
\begin{tabular}{lccc}
\toprule
\textbf{Метод} & \textbf{Accuracy (\%)} & \textbf{Число параметров (M)} & \textbf{Сложность (GFLOPs)} \\
\midrule
\textbf{Proposed CBAM-CNN (наш метод)} & \textbf{99.88 $\pm$ 0.037} & \textbf{0.15} & \textbf{$\sim$0.02} \\
Autoencoder-based CNN (Bao et al., 2024) & 93.10 & 44.00 & — \\
AC-DDPM + ResNet (Li et al., 2024) & 95.20 & 24.80 & — \\
WaferCap (Mishra et al., 2024) & 99.00 & — & — \\
MobileNetV2 & 97.56 & 2.23 & 0.004 \\
MobileNetV3\_Small & 97.41 & 1.53 & 0.001 \\
ShuffleNetV2\_x1.0 & 95.44 & 1.26 & 0.002 \\
EfficientNet-Lite0 & 95.72 & 4.31 & 0.005 \\
EfficientNetV2\_S & 98.33 & 20.19 & 0.113 \\
SwinTransformer\_T & 97.41 & 18.86 & 0.061 \\
Sparse feature learning (Altantawy \& Yakout, 2025) & 98.77 & 0.17 & — \\
Stacking ensemble model (Yang et al., 2025) & 98.18 & — & — \\
Soft voting approach (Lee et al., 2025) & 95.09 & — & — \\
\bottomrule
\end{tabular}
\end{table}

\noindent Модель CBAM–CNN обеспечивает точность 99.88\% на тестовом множестве, используя на 1–2 порядка меньше параметров, чем тяжелые трансформеры или глубокие CNN.

\section{Результаты анализа абляции и стресс-тестирования}

Для оценки вклада каждого отдельного компонента архитектуры и устойчивости модели были проведены абляционные исследования и тесты на устойчивость к шумам (на базе 3 случайных сидов):

\begin{table}[h!]
\centering
\caption{Анализ абляции и стресс-тестирование}
\label{tab:ablation}
\begin{tabular}{lccc}
\toprule
\textbf{Вариант конфигурации} & \textbf{Accuracy (\%)} & \textbf{Macro-F1 (\%)} & \textbf{ECE (\%)} \\
\midrule
\multicolumn{4}{l}{\textbf{(I) Абляция компонентов:}} \\
\textbf{Полная модель (CBAM + мультимасшт. + аугм.)} & \textbf{99.88 $\pm$ 0.05} & \textbf{99.83 $\pm$ 0.06} & \textbf{0.42 $\pm$ 0.08} \\
Без аугментации данных (CAE) & 98.72 $\pm$ 0.11 & 98.55 $\pm$ 0.13 & 0.96 $\pm$ 0.10 \\
Без CBAM (чистый сверточный бэкбон) & 98.11 $\pm$ 0.17 & 97.92 $\pm$ 0.19 & 1.21 $\pm$ 0.14 \\
Только канальное внимание (без пространственного) & 98.96 $\pm$ 0.10 & 98.80 $\pm$ 0.12 & 0.78 $\pm$ 0.09 \\
\midrule
\multicolumn{4}{l}{\textbf{(II) Стресс-тесты и дефицит данных:}} \\
С аддитивным шумом Гаусса на тесте ($\sigma = 0.05$) & 98.34 $\pm$ 0.15 & 98.19 $\pm$ 0.16 & 0.89 $\pm$ 0.11 \\
Обучение только на 25\% данных & 98.05 $\pm$ 0.18 & 97.84 $\pm$ 0.20 & 1.08 $\pm$ 0.13 \\
Обучение только на 10\% данных & 96.90 $\pm$ 0.25 & 96.41 $\pm$ 0.28 & 1.42 $\pm$ 0.17 \\
10\% данных + шум Гаусса ($\sigma = 0.05$) на тесте & 95.82 $\pm$ 0.31 & 95.21 $\pm$ 0.35 & 1.73 $\pm$ 0.22 \\
\bottomrule
\end{tabular}
\end{table}

\subsection*{Основные выводы абляции и тестов:}
\begin{itemize}
    \item Исключение блока внимания CBAM приводит к наиболее существенному падению качества классификации (снижение Macro-F1 на $\sim$1.9\%).
    \item Использование пространственного мультимасштабного внимания критически важно: удаление пространственного блока снижает Macro-F1 на $\sim$1.03\%.
    \item Модель демонстрирует высокую устойчивость к дефициту данных — даже при обучении всего на 10\% от всей выборки метрика Macro-F1 удерживается на высоком уровне 96.41\%.
\end{itemize}

\end{document}